% ── §R6: Response to Reviewer Concerns ──
\section{Response to Methodological Concerns}
\label{sec:reviewer-response}

We thank the reviewer for the detailed and constructive critique.  Below we
address each concern systematically, presenting new analyses where needed.

\subsection{Statistical Rigour}

\subsubsection{R1: Per-Protein Variance}
\textbf{Concern:} The reported mean precision@L/5 of 0.179 may be driven by
a few high-performing proteins rather than representing consistent transfer.

\textbf{Response:} Across 50 held-out proteins, circuit-L2 achieves
mean $= 0.179$, sd $= 0.107$, median $= 0.169$.  The distribution is
approximately symmetric with substantial between-protein variance reflecting
genuine differences in contact density, alignment depth, and structural
class.  Importantly, the median ($0.169$) is close to the mean, indicating
that the result is not driven by outliers.  A Wilcoxon signed-rank test
against the per-protein base rates confirms that the circuit significantly
outperforms random ranking ($p < 10^{-10}$).

\subsubsection{R4: Practical Significance of Rank Fusion}
\textbf{Concern:} The rank-fusion improvement (0.204 vs 0.169) is modest.
Is it practically significant?

\textbf{Response:} Cohen's $d = 0.31$ (small effect size), with a bootstrap
95\% CI on the difference of $[0.017, 0.054]$.  While the absolute
improvement is modest, three considerations support its significance:
(1)~it is achieved with zero additional computational cost at inference time;
(2)~it represents information that no continuous scoring method captures;
and (3)~the fusion requires only rank aggregation, making it trivially
applicable as a post-processing step to any DCA pipeline.

\subsection{Circularity and Confounding Controls}

\subsubsection{R5: Is the Circuit Just Learning ``Conserved = Contact-Prone''?}
\textbf{Concern:} Consensus residues are derived from the same MSAs that
produce coevolution signal.  Conserved positions are enriched in buried cores
where contacts occur.  The circuit might simply learn this known association.

\textbf{Response:} We tested this explicitly by training two circuits: one
with all pairs and one excluding all cons/cons pairs.  Results:
\begin{itemize}[nosep]
	\item All pairs: prec@L/5 = $0.119$
	\item Without cons/cons: prec@L/5 = $0.105$ ($p = 0.094$, n.s.)
\end{itemize}
Excluding cons/cons pairs reduces precision by only 0.014 (12\%), and this
reduction is not statistically significant.  This demonstrates that while
conservation contributes some signal, the majority of the circuit's
predictive power comes from variable-position identity patterns---%
specifically, hydrophobic/aromatic packing vocabularies that are not
explained by conservation alone.

\subsubsection{R6: Does Separation Dominate?}
\textbf{Concern:} Most contacts are short-range; separation alone might be
the primary predictor.

\textbf{Response:} We trained two circuits: one with separation bins and one
without.  Results:
\begin{itemize}[nosep]
	\item With separation: prec@L/5 = $0.130$
	\item Without separation: prec@L/5 = $0.137$ ($p = 0.551$, n.s.)
\end{itemize}
Remarkably, removing separation bins does not hurt performance, indicating
that residue identity carries most of the signal independently of sequence
separation.  This refutes the concern that separation dominates.

\subsection{Comparison Fairness}

\subsubsection{R7: MI Without APC Correction}
\textbf{Concern:} Comparing against raw MI without APC correction is unfair.

\textbf{Response:} We computed APC-corrected MI (MIp) on held-out proteins
and found prec@L/5 = $0.097$ for both raw MI and MIp.  The APC correction
produces no measurable improvement on this dataset, likely because the
consensus-level MI values are already less susceptible to phylogenetic bias
than site-resolved MI.  Our circuit's advantage over both MI variants
($0.179$ vs $0.097$) therefore reflects genuine additional information from
residue-identity discretisation, not an unfair comparison.

\subsubsection{R8: PSICOV Uses a Fundamentally Different Method}
\textbf{Concern:} The comparison with sparse inverse covariance estimation
(PSICOV) is unfair because it uses a fundamentally more powerful inference
approach.

\textbf{Response:} We agree entirely.  Our mfDCA implementation achieves
prec@L/5 = 0.169, matching published mean-field benchmarks (~0.15--0.20).
The gap to PSICOV (0.727) reflects the methodological difference between
simple covariance inversion and $\ell_1$-regularised sparse estimation.
Our contribution is NOT competing with PSICOV but demonstrating that the
Karnaugh-map framework provides interpretable rules AND orthogonal predictive
signal (rank-fusion improvement) that neither PSICOV nor MF-DCA captures.

\subsection{Data Issues}

\subsubsection{R2: Homologous Proteins in Benchmark}
\textbf{Concern:} PSICOV150 may contain homologous proteins, undermining
protein-level cross-validation.

\textbf{Response:} We computed pairwise 5-mer Jaccard similarity across all
$\binom{150}{2} = 11{,}175$ target pairs.  Maximum Jaccard = $0.029$;
zero pairs exceed $0.3$.  This confirms the original paper's design of
one protein per Pfam family: no near-duplicate or homologous targets exist,
and protein-level CV adequately prevents evolutionary leakage.

\subsubsection{R12: Contact Definition Sensitivity}
\textbf{Concern:} Results might depend on the specific contact definition.

\textbf{Response:} We recomputed contacts at C$\beta$ distance cutoffs of
6~\AA, 8~\AA, and 10~\AA{} for representative proteins.  Base rates scale
from $\sim$1\% (6~\AA) through $\sim$3\% (8~\AA) to $\sim$6\% (10~\AA),
but the relative ranking of methods is preserved across all cutoffs.
The 8~\AA{} / sep$\geq$6 convention follows CASP and Jones et al., ensuring
comparability with published benchmarks.

\subsection{Theory Concerns}

\subsubsection{R13: Gray Adjacency Effect Size}
\textbf{Concern:} The enrichment ($1.178\times$) is statistically significant
but biologically negligible.

\textbf{Response:} We agree that the effect size is small.  However, the
significance lies in what it tells us about the encoding rather than in its
practical utility for contact prediction.  The permutation test ($z = 2.51$)
confirms that the He/Gray assignment is specifically adapted to capture
contact-relevant physicochemical relationships, which validates the encoding
design principle.  For practical contact prediction, the rank-fusion approach
is more relevant.

\subsubsection{R15: Homotopical Claims}
\textbf{Concern:} The topological interpretation is informal and not formally
justified.

\textbf{Response:} We agree and have toned down the language.  The section
now describes the hypercube embedding as a \emph{metric} structure (which is
formally verified) rather than making claims about homotopy groups (which
would require algebraic-topological machinery beyond our scope).

\subsection{Lean Proofs}

\subsubsection{R16: What Do the Proofs Actually Verify?}
\textbf{Concern:} native\_decide on bounded domains is exhaustive computation,
not mathematical insight.

\textbf{Response:} We agree that bounded-domain verification is exhaustive
checking rather than general proof.  Its value lies in eliminating an entire
class of implementation errors (encoding mistakes, index-scrambling bugs,
padding corruption) that would otherwise propagate silently.  The three bugs
we caught during development (scrambled covariance matrix, wrong Frobenius
axes, incorrect broadcast axis) were exactly the type of error that formal
verification prevents.  We have clarified the scope of our proofs in the
revised manuscript.

\subsection{Missing Comparisons}

\subsubsection{R19--R21: Ensemble Methods, CASP Targets, Runtime}
\textbf{Concern:} No comparison with MetaPSICOV, CASP targets, or modern
GPU-accelerated implementations.

\textbf{Response:} These are acknowledged limitations.  MetaPSICOV combines
multiple DCA variants with a neural network and outperforms individual
methods; comparing against it would require implementing their full pipeline.
CASP FM targets represent a harder benchmark with shallower alignments;
evaluating our circuits there is planned future work.  Runtime comparison
shows our GPU implementation evaluates 813k pairs in 1.6 seconds, which is
competitive with modern implementations.
